\cvsection{Expériences professionnelles}

\begin{cventries}
  \cventry
    {Stagiaire ingénieur de recherche en IA}
    {INRIA}
    {Bordeaux, France}
    {Mars 2026 -- Septembre 2026}
    {
      \begin{cvitems}
        \item {Conception en autonomie, avec le suivi de mes encadrants, d’une application IA fonctionnelle de bout en bout pour aider les chercheurs à trouver des articles scientifiques.}
        \item {Développement d’un backend Python/FastAPI avec API REST documentées avec OpenAPI et persistance des documents et métadonnées dans MongoDB.}
        \item {Développement d’une interface React/Next.js pour interagir avec l’agent et ses outils, avec sélection par l’utilisateur des messages à inclure dans le contexte du LLM.}
        \item {Construction d’un agent multi-étapes sans framework d’orchestration, avec API OpenAI, prompts adaptés aux outils et gestion de l’état et de la mémoire.}
        \item {Intégration d’un pipeline RAG avec chunking, embeddings et recherche sémantique dans une base vectorielle ChromaDB.}
        \item {Automatisation de la collecte de métadonnées d’articles scientifiques en ligne et transmission des résultats de tâches cron par webhooks.}
        \item {Conteneurisation des composants avec Docker et mise en place d’une CI/CD sur GitHub.}
        \item {Gestion d’une application utilisant des LLM installés localement sur un Mac Studio et utilisation de vLLM pour servir des modèles sur HPC avec GPU NVIDIA.}
        \item {Comparaison de Qwen, Gemma, GLM et de plusieurs quantifications avec Slurm selon la qualité, la latence, le TTFT et le contexte, avec rédaction en cours d’un article scientifique.}
        \item {Évaluation d’OpenClaw dans un environnement Docker isolé et analyse des risques liés à l’accès autonome de l’agent aux outils et aux fichiers.}
      \end{cvitems}
    }

  \cventry
    {Stagiaire de recherche en apprentissage par renforcement}
    {Université Karunya}
    {Coimbatore, Inde}
    {Juin 2025 -- Sept. 2025}
    {
      \begin{cvitems}
        \item {Développement d’un modèle d’apprentissage par renforcement et de son environnement d’entraînement pour la navigation autonome d’un robot hospitalier.}
        \item {Déploiement et optimisation du logiciel et du modèle sur Raspberry Pi pour des tests en conditions réelles.}
      \end{cvitems}
    }
\end{cventries}
